\chapter{Financing: Repo, Securities Lending and Prime Brokerage}\label{ch:m1:financing}

A fund owns five billion dollars of stock and has one billion of its own
money. The other four are borrowed, and the loan is renewed every night by
lenders who can decline to renew it. On an ordinary day nobody thinks about
this. On the day the stocks fall four percent, a fifth of the fund's money
is gone, its lenders want more collateral by the afternoon, and the only way
to find it is to sell \$800 million of the same stocks into the same falling
market. Almost every position in professional finance is carried with
somebody else's money. This chapter is about where that money comes from,
what it costs, the three contracts through which it is lent --- and what
happens when it leaves.

\section{Leverage}

\begin{definition}[Leverage]\label{def:m1:financing:leverage}
The \emph{leverage}\index{leverage} of a portfolio is the ratio $L = A/E$ of
the assets it holds to the equity (own funds) behind them. The remainder,
$A - E$, is financed.
\end{definition}

\begin{proposition}[Leverage arithmetic]\label{prop:m1:financing:roe}
Let the assets return $r_A$ over a period and the financing cost $r_f$. Then
\begin{enumerate}
\item the return on equity is $r_E = L\,r_A - (L-1)\,r_f$;
\item the equity is wiped out when the assets fall by $1/L$;
\item to restore the \omterm{def:m1:financing:leverage}{leverage} ratio after the assets fall by a fraction $x$
(with $x < 1/L$), the portfolio must sell assets worth $(L-1)\,x\,A$.
\end{enumerate}
\end{proposition}
\begin{proof}
(1) The equity earns $A r_A - (A-E) r_f$; divide by $E$. (2) The loss $xA$
equals $E$ when $x = E/A$. (3) After the fall assets are $A(1-x)$ and equity
is $E - xA$. Assets consistent with \omterm{def:m1:financing:leverage}{leverage} $L$ are $L(E - xA) = A - LxA$.
The excess to sell is $A(1-x) - A + LxA = (L-1)xA$.
\end{proof}

\begin{example}[The fund of the opening paragraph]\label{ex:m1:financing:fund}
$A = \$5$ billion, $E = \$1$ billion, $L = 5$. With $r_f = 4\%$, a $+10\%$
year on the assets is $+34\%$ on the equity; a $-10\%$ year is $-66\%$. A
fall of 20\% ends the fund. After a fall of 4\% the fund has lost \$200
million, a fifth of its equity, and must sell $4 \times 0.04 \times 5 =
\$0.8$ billion to be five times leveraged again: \emph{four dollars of
selling for each dollar lost}.
\end{example}

\begin{omfigure}
\begin{tikzpicture}
\begin{axis}[width=10.5cm,height=5cm,
  xlabel={return of the assets (\%)},ylabel={return on equity (\%)},
  tick label style={font=\small},label style={font=\small},
  xmin=-20,xmax=20,ymin=-110,ymax=140,grid=major,grid style={gray!25},
  legend columns=3,legend style={font=\footnotesize,at={(0.5,-0.32)},anchor=north,draw=none,fill=none,/tikz/every even column/.append style={column sep=8pt}},
  table/col sep=comma]
\addplot[omMeth,very thick] table[x=asset_ret_pct,y=l1]{figdata/markets-1/06-financing/roe.csv};
\addplot[omDef,very thick] table[x=asset_ret_pct,y=l3]{figdata/markets-1/06-financing/roe.csv};
\addplot[omThm,very thick] table[x=asset_ret_pct,y=l8]{figdata/markets-1/06-financing/roe.csv};
\addplot[black,dashed] coordinates {(-20,-100) (20,-100)};
\legend{$L=1$,$L=3$,$L=8$}
\end{axis}
\end{tikzpicture}

\omcaption{Return on equity against the return of the assets, with financing
at 4\%. The dashed line is total loss: at $L = 8$ it is reached when the
assets fall 9\% --- less than $1/L = 12.5\%$, because a year of financing
cost is owed as well. Data: \cref{prop:m1:financing:roe}, computed by the chapter's
script.}
\end{omfigure}

Who lends the $A - E$? Three contracts do almost all of the work: the repo,
the securities loan, and the swap. All three are \emph{secured}: the lender
holds collateral worth more than the loan, and that excess is the quantity
to watch.

\section{Repo}

\begin{definition}[Repurchase agreement and haircut]\label{def:m1:financing:repo}
In a \emph{repurchase agreement}\index{repurchase agreement} (repo) one party
sells a security for cash and agrees at the same time to buy it back at a
fixed later date at a fixed higher price. Economically it is a cash loan
secured by the security; the price difference is the interest, quoted as the
\emph{repo rate}. The \emph{haircut}\index{haircut} $h$ is the fraction by
which the cash lent falls short of the security's market value: a security
worth 100 raises $100(1-h)$.
\end{definition}

\begin{omfigure}
\begin{tikzpicture}[font=\small,
  box/.style={draw=omDef,rounded corners=2pt,minimum width=3cm,minimum height=1cm,align=center,fill=omDef!4},
  flow/.style={-{Stealth[length=2mm]},thick}]
\node[box] (b1) at (0,1.5) {Cash borrower\\(security owner)};
\node[box] (l1) at (9,1.5) {Cash lender};
\draw[flow,omDef] ([yshift=5pt]b1.east) -- node[above,font=\footnotesize]{bond worth 100} ([yshift=5pt]l1.west);
\draw[flow,omThm] ([yshift=-5pt]l1.west) -- node[below,font=\footnotesize]{cash 98 \quad (haircut 2\%)} ([yshift=-5pt]b1.east);
\node[font=\footnotesize\bfseries,anchor=east] at (-2.0,1.5) {Today};
\node[box] (b2) at (0,-1.1) {Cash borrower};
\node[box] (l2) at (9,-1.1) {Cash lender};
\draw[flow,omThm] ([yshift=5pt]b2.east) -- node[above,font=\footnotesize]{cash $98\,(1 + r\,\tfrac{n}{360})$} ([yshift=5pt]l2.west);
\draw[flow,omDef] ([yshift=-5pt]l2.west) -- node[below,font=\footnotesize]{the same bond} ([yshift=-5pt]b2.east);
\node[font=\footnotesize\bfseries,anchor=east] at (-2.0,-1.1) {In $n$ days};
\end{tikzpicture}

\omcaption{The two legs of a repo. If the borrower fails to repurchase, the
lender owns a bond worth 100 against a claim of 98: the \omterm{def:m1:financing:repo}{haircut} is its
protection against a fall in the bond's price while it sells.}
\end{omfigure}

\begin{proposition}[Haircuts bound leverage]\label{prop:m1:financing:haircut}
A portfolio financed entirely by repo at \omterm{def:m1:financing:repo}{haircut} $h$ can reach a \omterm{def:m1:financing:leverage}{leverage} of
at most $1/h$. If the \omterm{def:m1:financing:repo}{haircut} rises from $h$ to $h'$ with no change in
prices, a portfolio at maximum \omterm{def:m1:financing:leverage}{leverage} must sell a fraction $1 - h/h'$ of
its assets.
\end{proposition}
\begin{proof}
Each unit of assets raises $1-h$ of cash, so the equity needed is $hA$ and
$L = A/E \le 1/h$. With equity unchanged at $E = hA$, the new maximum is
$E/h' = (h/h')A$.
\end{proof}

\begin{example}[Two percent to four percent]\label{ex:m1:financing:doubling}
Government bonds financed at a 2\% \omterm{def:m1:financing:repo}{haircut} allow $L = 50$: this is how
relative-value funds earn a living from price differences of a few basis
points (One Quant Book 9). If lenders move the \omterm{def:m1:financing:repo}{haircut} to 4\%, a fund at the
limit must sell \emph{half} its assets although no price has moved. \omterm{def:m1:financing:repo}{Haircuts}
are set by lenders, rise when volatility rises, and are the channel through
which a funding problem becomes a market problem.
\end{example}

\section{Securities lending and the short sale}

\begin{definition}[Short sale]\label{def:m1:financing:short}
A \emph{short sale}\index{short sale} is the sale of a security the seller
does not own. To deliver at settlement the seller borrows the security, and
later buys it back to return it: it profits if the price has fallen.
\end{definition}

\begin{definition}[Securities lending, rebate rate, locate]\label{def:m1:financing:seclending}
In \emph{securities lending}\index{securities lending} an owner transfers a
security to a borrower against collateral, usually cash worth a little more
than the security, and receives an equivalent security back on demand. The
lender invests the cash and returns part of the interest to the borrower at
the \emph{rebate rate}\index{rebate rate}: the difference between the market
interest rate and the rebate is the \emph{borrow fee}\index{borrow fee}, the true price of the
loan. A \emph{locate}\index{locate} is a broker's confirmation, obtained
before a \omterm{def:m1:financing:short}{short sale}, that the security can be borrowed in time for
settlement; US rules require one (\cref{ch:m1:stock-loan-and-short-selling}).
\end{definition}

\begin{omfigure}
\begin{tikzpicture}[font=\small,
  box/.style={draw=omDef,rounded corners=2pt,minimum width=2.7cm,minimum height=0.95cm,align=center,fill=omDef!4},
  flow/.style={-{Stealth[length=2mm]},thick}]
\node[box] (len) at (0,0) {Lender\\(a pension fund)};
\node[box] (pb) at (5.2,0) {Prime broker};
\node[box] (ss) at (10.4,0) {Short seller};
\node[box] (mk) at (10.4,-2.6) {Market};
\draw[flow,omDef] ([yshift=6pt]len.east) -- node[above,font=\footnotesize]{shares} ([yshift=6pt]pb.west);
\draw[flow,omThm] ([yshift=-6pt]pb.west) -- node[below,font=\footnotesize]{collateral} ([yshift=-6pt]len.east);
\draw[flow,omDef] ([yshift=6pt]pb.east) -- node[above,font=\footnotesize]{shares} ([yshift=6pt]ss.west);
\draw[flow,omThm] ([yshift=-6pt]ss.west) -- node[below,font=\footnotesize]{proceeds} ([yshift=-6pt]pb.east);
\draw[flow,omDef] ([xshift=-6pt]ss.south) -- node[left,font=\footnotesize]{shares sold} ([xshift=-6pt]mk.north);
\draw[flow,omThm] ([xshift=6pt]mk.north) -- node[right,font=\footnotesize]{proceeds} ([xshift=6pt]ss.south);
\end{tikzpicture}

\omcaption{A \omterm{def:m1:financing:short}{short sale}. Shares (blue) travel from a long-term owner to the
market; cash (red) travels back and stops as collateral with the lender,
who pays interest on it at the \omterm{def:m1:financing:seclending}{rebate rate}. The short seller never holds the
proceeds.}
\end{omfigure}

\begin{example}[The carry of a short]\label{ex:m1:financing:carry}
Interest rates are 4\%. A widely held large stock lends at a fee of 0.25\%:
the rebate is 3.75\%, and a short seller \emph{earns} 3.75\% a year on the
proceeds. A stock in heavy demand lends at a fee of 15\%: the rebate is
$-11\%$, and the short seller pays 11\% a year for the right to be short. In
both cases the short also pays the lender every dividend the stock
distributes.
\end{example}

\section{Prime brokerage and margin}

\begin{definition}[Prime broker and rehypothecation]\label{def:m1:financing:pb}
A \emph{prime broker}\index{prime broker} is the bank that finances a \omterm{def:m1:what-a-trading-firm-does:hedge-fund}{hedge
fund}'s securities portfolio: it lends cash against longs, lends securities
for shorts, settles and holds the portfolio, and reports on it.
\emph{Rehypothecation}\index{rehypothecation} is the prime broker's re-use of
a client's securities as collateral for its own borrowing, which is how it
funds the loans it makes.
\end{definition}

\begin{definition}[Margin call]\label{def:m1:financing:margincall}
A \emph{margin call}\index{margin call} is a lender's demand that a borrower
restore the agreed excess of collateral over loan, by delivering cash or
securities or by reducing positions, usually within one day. If it is not
met the lender may sell the collateral.
\end{definition}

\begin{dated}{2026-09}{US margin rules in three numbers}\label{dat:m1:financing:rules}
The Federal Reserve's Regulation T sets the \emph{initial} margin on a
purchase of listed stock in a margin account at 50\% (and 150\% for a \omterm{def:m1:financing:short}{short
sale}, of which the sale proceeds provide 100\%). FINRA Rule 4210 sets the
\emph{maintenance} margin at 25\% of market value for long listed equities.
Under SEC Rule 15c3-3 a broker may use a customer's securities up to 140\% of
the customer's debit balance; securities beyond that amount must be kept in
the broker's possession or control. \omterm{def:m1:financing:pb}{Prime brokers}' own house requirements
and portfolio-margin regimes sit on top of these.
\end{dated}

\begin{proposition}[Where the margin call is]\label{prop:m1:financing:call}
A long position bought at $P_0$ with \omterm{def:m1:clearing-and-settlement:margin}{initial margin} ratio $m_0$ (equity over
value) receives a maintenance call at maintenance ratio $m$ when the price
falls below
\[
P^\star = P_0\,\frac{1-m_0}{1-m}.
\]
\end{proposition}
\begin{proof}
The loan is $P_0(1-m_0)$ and does not change with the price. At price $P$
the equity ratio is $(P - P_0(1-m_0))/P$; set it equal to $m$.
\end{proof}

\begin{example}[Fifty and twenty-five]\label{ex:m1:financing:regt}
With $m_0 = 50\%$ and $m = 25\%$, $P^\star = \tfrac23 P_0$: the call comes
after a fall of a third. A professional client margined by its \omterm{def:m1:financing:pb}{prime broker}
at $m_0 = 15\%$ and $m = 10\%$ is called at $0.944\,P_0$, after a fall of
5.6\%.
\end{example}

\section{Synthetic financing}

\begin{definition}[Total return swap]\label{def:m1:financing:trs}
In an equity \emph{total return swap}\index{total return swap} the dealer
pays the client the total return (price change and dividends) of a stock on
a notional amount, and the client pays a financing rate plus a spread on the
same notional. The client has the economics of owning the stock financed by
the dealer, who holds the actual shares as its hedge; the client posts
margin agreed in the contract.
\end{definition}

The swap is financing in another wrapper, and the wrapper matters: the
client appears on no shareholder register; the margin is whatever the two
parties negotiated, not what Regulation T prescribes; and a client using
several dealers shows each of them only its own slice.

\begin{example}[Archegos]\label{ex:m1:financing:archegos}
In March 2021 a family office, Archegos Capital Management, defaulted on
\omterm{def:m1:financing:margincall}{margin calls} from its dealers. It had built concentrated positions in a
handful of stocks through \omterm{def:m1:financing:trs}{total return swaps} with several banks, positions
which the SEC's complaint later put at \$36 billion. One of those banks,
Credit Suisse, lost close to \$5.5 billion. The report its board
commissioned found that the bank had agreed to a swap margin of 7.5\% ---
\omterm{def:m1:financing:leverage}{leverage} above thirteen --- that the margin was \emph{static}, fixed on the
price at which each swap was opened, so that as the stocks rose the average
margin held fell to 6.9\% of current value, and that a move to dynamic
margining had not been given priority. When the stocks fell, every dealer
held the same shares as its hedge, and each had to sell them into the
others' selling.
\end{example}

\section{When financing runs}

Combine \cref{prop:m1:financing:roe}(3) with the fact that large sales move
prices (\cref{met:m1:the-sell-side:sqrt}) and the mechanism of every
leveraged crisis appears: a fall forces sales, sales cause a further fall,
which forces more sales.

\begin{omfigure}
\begin{tikzpicture}
\begin{axis}[width=10.5cm,height=4.8cm,
  xlabel={round of forced selling},ylabel={cumulative price fall (\%)},
  tick label style={font=\small},label style={font=\small},
  xmin=1,xmax=12,ymin=4,ymax=11,grid=major,grid style={gray!25},
  legend columns=3,legend style={font=\footnotesize,at={(0.5,-0.34)},anchor=north,draw=none,fill=none,/tikz/every even column/.append style={column sep=8pt}},
  table/col sep=comma]
\addplot[omMeth,very thick,mark=*,mark size=1.2pt] table[x=round,y=l3]{figdata/markets-1/06-financing/spiral.csv};
\addplot[omDef,very thick,mark=*,mark size=1.2pt] table[x=round,y=l5]{figdata/markets-1/06-financing/spiral.csv};
\addplot[omThm,very thick,mark=*,mark size=1.2pt] table[x=round,y=l8]{figdata/markets-1/06-financing/spiral.csv};
\legend{$L=3$,$L=5$,$L=8$}
\end{axis}
\end{tikzpicture}

\omcaption{A 5\% shock, then rounds of selling to restore \omterm{def:m1:financing:leverage}{leverage}, each sale
moving the price by 10\% per unit of assets sold (an illustrative impact).
At $L = 8$ the market ends twice as far down as the shock that started it.
Data: the chapter's script.}\label{fig:m1:financing:spiral}
\end{omfigure}

\begin{remark}[Three accelerants]\label{rem:m1:financing:accelerants}
The toy spiral of \cref{fig:m1:financing:spiral} converges. Real ones are
made worse by three things it leaves out. Lenders raise \omterm{def:m1:financing:repo}{haircuts}
(\cref{prop:m1:financing:haircut}) exactly when prices fall. Funds holding
the same positions are hit together, so the ``market'' absorbing the sales is
itself selling. And lenders who fear for a borrower stop renewing overnight
loans altogether, turning a \omterm{def:m1:financing:margincall}{margin call} into a run. The quantitative-fund
losses of August 2007, the repo run of 2008 and the 2021 episode above are
the three cases this series returns to.
\end{remark}

\section{Tutorial: what a long--short book costs to carry}\label{tut:m1:financing:carry}

\begin{tutorial}
\textbf{Goal.} Compute the annual financing cost of a long--short equity
book, first with a cash prime-brokerage account, then with swaps.
\textbf{End state:} two numbers, \$2.2 million and \$2.7 million, and the
reason for the difference.
\begin{tutsteps}
\item \textbf{The book.} \$300 million long, \$200 million short, \$100
million of equity. Illustrative terms: interest rate 4\%; the \omterm{def:m1:financing:pb}{prime broker}
charges 0.50\% over on debit balances and pays 0.30\% under on short
proceeds; the average \omterm{def:m1:financing:seclending}{borrow fee} is 0.30\%.
\item \textbf{Cash account.} The fund borrows $300 - 100 = \$200$ million.
\omcode{code/markets-1/06-financing/python/financing_demo.py}{60}{64}{Cash prime brokerage: pay on the debit, receive on the short proceeds.}\label{lst:m1:financing:cash}
Cost: $200 \times 4.5\% - 200 \times 3.4\% = \$2.2$ million a year.
\item \textbf{Swaps.} The dealer finances the whole long notional and the
fund's equity sits in cash earning the interest rate.
\omcode{code/markets-1/06-financing/python/financing_demo.py}{67}{72}{The same book on swap.}\label{lst:m1:financing:swap}
Cost: $300 \times 4.5\% - 200 \times 3.4\% - 100 \times 4\% = \$2.7$ million.
\item \textbf{Explain the gap.} \$0.5 million is the 0.50\% spread paid on
the \$100 million of longs that the cash account funded with the fund's own
money. Synthetic financing charges the spread on everything.
\end{tutsteps}
\textbf{What to change next.} Put a fifth of the short book in stocks
lending at a 12\% fee. Then let the interest rate fall to zero: which of the
two costs changes, and why does neither change much?
\end{tutorial}

\section{Build: the financing calculator}\label{bld:m1:financing:calc}

\begin{build}
\bfield{Purpose} Each night the miniature firm accrues what its positions
cost to carry and posts it to the ledger of
\cref{ch:m1:what-a-trading-firm-does} under \texttt{financing}.
\bfield{Interface} \texttt{FinancingTerms(debit\_spread, credit\_spread,
day\_count=360)}, a per-symbol table of \omterm{def:m1:financing:seclending}{borrow fees}, and
\texttt{accrue(date, positions, prices, equity, rate)} returning one signed
amount per line: debit interest, short-proceeds interest, \omterm{def:m1:financing:seclending}{borrow fees}.
\bfield{Rules} Accrue on settled, not traded, positions. Interest is simple,
actual days over the day count, so a Friday accrues three days. \omterm{def:m1:financing:seclending}{Borrow fees}
accrue on the value of the short marked at the previous close. A missing
\omterm{def:m1:financing:seclending}{borrow fee} for a short position is an error, not a zero.
\bfield{Acceptance tests} \texttt{code/firm/financing/tests/}: reproduces the
tutorial's \$2.2 million over a 360-day year; a weekend accrues three days;
a missing fee raises.
\bfield{Stretch} Add \omterm{def:m1:financing:repo}{haircuts} per asset class and a \texttt{max\_financing}
query: how much more can the firm buy today?
\end{build}

\begin{omsources}
\item Credit Suisse Group Special Committee of the Board of Directors,
\emph{Report on Archegos Capital Management}, 29 July 2021.
\item US Securities and Exchange Commission, \emph{SEC Charges Archegos and
its Founder with Massive Market Manipulation Scheme}, press release 2022-70,
27 April 2022.
\item 12 CFR 220.12 (Regulation T, supplement: margin requirements); FINRA
Rule 4210; SEA Rule 15c3-3 (FINRA interpretations handbook).
\item M.~Brunnermeier and L.~Pedersen, ``Market liquidity and funding
liquidity'', \emph{Review of Financial Studies} 22 (2009).
\item G.~Gorton and A.~Metrick, ``Securitized banking and the run on repo'',
\emph{Journal of Financial Economics} 104 (2012).
\end{omsources}

\section{Exercises}

\begin{exercise}[$\star$]\label{exo:m1:financing:1}
A fund has equity of \$250 million and assets of \$1.5 billion. Give its
\omterm{def:m1:financing:leverage}{leverage}, the fall in its assets that wipes it out, and its return on equity
if the assets return $+6\%$ with financing at 4\%.
\end{exercise}

\begin{exercise}[$\star$]\label{exo:m1:financing:2}
A bond worth \$50 million is financed in repo at a \omterm{def:m1:financing:repo}{haircut} of 3\% and a \omterm{def:m1:financing:repo}{repo
rate} of 4.2\% for 7 days (actual/360). How much cash is raised and how much
interest is paid?
\end{exercise}

\begin{exercise}[$\star$]\label{exo:m1:financing:3}
Interest rates are 4\%. A short seller is short \$20 million of a stock with
a \omterm{def:m1:financing:seclending}{borrow fee} of 6\%. What does the position earn or cost per year in
financing, before dividends?
\end{exercise}

\begin{exercise}[$\star\star$]\label{exo:m1:financing:4}
A stock is bought at \$80 with 40\% \omterm{def:m1:clearing-and-settlement:margin}{initial margin}; the maintenance ratio is
30\%. At what price is the \omterm{def:m1:financing:margincall}{margin call}? The stock falls to \$60: how much
cash per share restores the \emph{initial} ratio?
\end{exercise}

\begin{exercise}[$\star\star$]\label{exo:m1:financing:5}
A fund with $L = 6$ and assets of \$3 billion loses 5\% on its assets.
Compute its new \omterm{def:m1:financing:leverage}{leverage}, and the sale needed to return to $L = 6$. If its
lenders now allow only $L = 4$, how much must it sell?
\end{exercise}

\begin{exercise}[$\star\star$]\label{exo:m1:financing:6}
A relative-value fund runs at the maximum \omterm{def:m1:financing:leverage}{leverage} allowed by a 2.5\%
\omterm{def:m1:financing:repo}{haircut}. Lenders raise the \omterm{def:m1:financing:repo}{haircut} to 4\%. What fraction of the book must be
sold? The trade earns 0.12\% a year on assets above its financing cost:
give the return on equity before and after.
\end{exercise}

\begin{exercise}[$\star\star\star$]\label{exo:m1:financing:7}
\emph{Coding.} With \texttt{spiral}, find by bisection the impact per unit
sold above which a 5\% shock wipes out a fund with $L = 8$ within twelve
rounds. Compare with the value $1/(L-1)$ and explain the relationship.
\end{exercise}

\begin{exercise}[$\star\star\star$]\label{exo:m1:financing:8}
\emph{Find the flaw.} A dealer's risk report says of a swap client: ``Margin
held is 7.5\% of notional; the largest one-day move of the portfolio in the
last five years was 6\%; the exposure is fully covered.'' The margin is
static and the client's stocks have doubled since the swaps were opened.
Identify three independent errors.
\end{exercise}

\section{Problem: A Family Office on Swap}

\begin{problem}[{Weekend problem --- leverage through five dealers}]\label{pb:m1:financing:1}
A family office has \$4 billion of equity. Through \omterm{def:m1:financing:trs}{total return swaps} with
five dealers, in equal shares, it holds long positions in eight stocks with
a current value of \$28 billion. Each dealer took a static margin of 10\% of
the price at which its swaps were opened; the stocks have risen 60\% on
average since.

\medskip\noindent\textbf{Part I --- How leveraged is it?}
\begin{enumerate}
\item Give the \omterm{def:m1:financing:leverage}{leverage} on current values and the fall that wipes out the
equity.
\item What was the notional when the swaps were opened, and how much margin
do the dealers hold in total?
\item Express that margin as a percentage of current value.
\item Where is the rest of the family office's \$4 billion? Why does it
matter to the dealers that they cannot see it?
\end{enumerate}

\medskip\noindent\textbf{Part II --- The carry.}
\begin{enumerate}[resume]
\item The swaps cost the interest rate (4\%) plus 0.45\%. Give the annual
financing cost in dollars.
\item The family office's cash earns 4\%. Give its net annual carry cost and
express it as a percentage of equity.
\item What annual return on the stocks does it need to break even?
\end{enumerate}

\medskip\noindent\textbf{Part III --- The fall.}
\begin{enumerate}[resume]
\item The stocks fall 10\% in a day. Give the loss, the remaining equity and
the new \omterm{def:m1:financing:leverage}{leverage}.
\item Each dealer calls for \omterm{def:m1:clearing-and-settlement:margin}{variation margin} equal to the loss on its slice.
How much cash is called in total? The family office pays out all its free
cash, pro rata. What does each dealer receive, and what is each still owed?
\item The next day the stocks fall another 8\%. Repeat question 8.
\item The family office is in default. Give the value of one dealer's hedge
shares after both falls, its total loss on the swaps over the two days, and
what it holds against that loss (margin plus cash received).
\item A dealer sells its shares over the following days at an average 12\%
below the second day's close. What is its final loss?
\item What would it have lost had it sold immediately at 3\% below that
close?
\end{enumerate}

\medskip\noindent\textbf{Part IV --- What should have been different.}
\begin{enumerate}[resume]
\item With dynamic margin of 10\% of \emph{current} value, how much margin
would the dealers have held before the fall?
\item Each position is about five days of the stock's trading volume for one
dealer alone. With a daily volatility of 3\%, estimate the 95\% discount of
\cref{prop:m1:the-sell-side:discount} for a dealer liquidating alone at 10\%
participation, and compare with the margin held.
\item Why is that estimate still too low in this situation?
\item What single piece of information, had the dealers shared it, would
have changed their margin?
\item Why did the first dealer to sell lose least?
\item State the \emph{named result}: the fall in the stocks that wipes out
the family office's equity, in percent.
\item In one sentence: what is the difference between \omterm{def:m1:financing:leverage}{leverage} and margin?
\end{enumerate}
\end{problem}

\section{Interview questions}

\begin{interviewq}[$\star$ \iqroles{trader, researcher, bank}]\label{iq:m1:financing:1}
Walk me through the cash flows of a \omterm{def:m1:financing:short}{short sale}, from the \omterm{def:m1:financing:seclending}{locate} to the
buy-back.
\end{interviewq}

\begin{interviewq}[$\star$ \iqroles{bank, trader}]\label{iq:m1:financing:2}
What is a repo, and why is it considered safe for the cash lender?
\end{interviewq}

\begin{interviewq}[$\star\star$ \iqroles{researcher, trader}]\label{iq:m1:financing:3}
A fund is leveraged four times and loses 5\% on its assets. How much must it
sell to keep its \omterm{def:m1:financing:leverage}{leverage} constant? Do it in your head.
\end{interviewq}

\begin{interviewq}[$\star\star$ \iqroles{researcher}]\label{iq:m1:financing:4}
Your backtest of a long--short strategy shows 9\% a year. It ignores
financing. List what you must subtract and give orders of magnitude.
\end{interviewq}

\begin{interviewq}[$\star\star$ \iqroles{bank}]\label{iq:m1:financing:5}
Why would a client choose a \omterm{def:m1:financing:trs}{total return swap} over buying the shares through
its \omterm{def:m1:financing:pb}{prime broker}? Why would the bank prefer one or the other?
\end{interviewq}

\begin{interviewq}[$\star\star\star$ \iqroles{bank, researcher}]\label{iq:m1:financing:6}
You are the risk manager of a \omterm{def:m1:financing:pb}{prime broker}. A client's portfolio is
concentrated in five stocks, in each of which it holds several days of
volume, and you suspect it has the same positions at other banks. How do you
set its margin?
\end{interviewq}
